\maketitle%与上面的\title对应
\begin{abstract}%摘要

	干燥是决定中药材成品品质的关键工序之一，热风烘干通过调控烘房温湿环境实现药材脱水，其工艺参数的合理选取直接决定干燥效率、能耗与成品品质。本文以近似圆柱形的药材为对象，建立传热—传质耦合的偏微分方程模型，对预热平衡与恒温干燥全过程进行数值仿真，系统研究了药材内部温度与水分浓度的时空演化规律。

	\textbf{对于问题一}，将药材视为长圆柱，建立一维径向非稳态热传导与水分扩散耦合模型（附录 2 常物性），采用 Crank--Nicolson 隐式有限差分法离散求解。结果表明：预热平衡阶段热量由表面向中心逐渐渗透，$1800$ s 时表面温度升至 $36.79$ °C、中心仅 $33.58$ °C；水分仅在表层下降，表面含水率由 $2.55$ 降至 $1.510$ kg/kg，中心基本不变。

	\textbf{对于问题二}，引入随含水率与温度变化的物性经验公式（附录 3），建立覆盖整个烘干过程的变物性模型，预热段采用附件 1 实测烘房温湿度，恒温段取 $50$ °C、$0.05$ kg/kg。结果表明：$3$ h 内药材温度趋于均匀并逼近烘房温度（$49.85\sim49.97$ °C），水分沿径向由表及里持续下降，中心与表面含水率分别降至 $1.766$、$1.008$ kg/kg。

	\textbf{对于问题三}，以药材各处水分浓度均低于 $0.15$ kg/kg 为烘干判据，由变物性模型求解得到\textbf{烘干结束时间约为 $57.45$ h（约 2.39 天）}，并给出每隔 6 h 的水分浓度分布。

	\textbf{对于问题四}，考虑药材失水导致的尺寸收缩，由附件 2 半径数据经保单调三次样条插值构造移动边界，结合附录 4 物性公式求解，得到\textbf{烘干时长为 $52.63$ h（约 2.19 天）}，收缩使扩散路径缩短，较定半径情形略有加快。

	进一步通过网格收敛性检验验证数值解二阶收敛，并对干燥时间开展单因素灵敏度与 SHAP 贡献度分析，结果表明空气湿度、对流传质系数与初始含水率是影响干燥时间的主导因素。本文模型物理机理清晰、数值结果稳定可靠，可为药材烘干工艺参数优化提供理论依据与决策参考。

	\keywords{药材烘干；传热传质；Crank--Nicolson 有限差分；干燥时间；移动边界}
\end{abstract}
